\documentclass[10pt,a4paper]{amsart}
\usepackage{amsmath,amssymb,amsthm,graphicx}
% Retained minimal preamble for the Stacks Project chapter more-algebra.tex (More on Algebra).
% The chapter file pinned-original-source.tex carries only \input{preamble} on its first line, so no
% preamble text lives inside the chapter itself.  The Stacks Project uses one single shared preamble for
% every chapter; the snapshot stages it identically as sources/stacks-chow/preamble.tex and
% sources/stacks-etale-cohomology/preamble.tex.  The block below is copied mechanically, line for line,
% from sources/stacks-chow/preamble.tex lines 14-18 (the Xy-pic setup) and 154-217 (the native theorem
% declarations together with the mathematical macros).  Nothing was added, removed or redefined; the
% wrapper class and the amsmath/amssymb/amsthm/graphicx base packages are supplied by the assembler.
% Chapter: more-algebra.tex -- More on Algebra.
% For commutative diagrams we use Xy-pic
\usepackage[all]{xy}

% We use 2cell for 2-commutative diagrams.
\xyoption{2cell}

% Theorem environments.
%
\theoremstyle{plain}
\newtheorem{theorem}[subsection]{Theorem}
\newtheorem{proposition}[subsection]{Proposition}
\newtheorem{lemma}[subsection]{Lemma}

\theoremstyle{definition}
\newtheorem{definition}[subsection]{Definition}
\newtheorem{example}[subsection]{Example}
\newtheorem{exercise}[subsection]{Exercise}
\newtheorem{situation}[subsection]{Situation}

\theoremstyle{remark}
\newtheorem{remark}[subsection]{Remark}
\newtheorem{remarks}[subsection]{Remarks}

\numberwithin{equation}{subsection}

% Macros
%
\def\lim{\mathop{\mathrm{lim}}\nolimits}
\def\colim{\mathop{\mathrm{colim}}\nolimits}
\def\Spec{\mathop{\mathrm{Spec}}}
\def\Hom{\mathop{\mathrm{Hom}}\nolimits}
\def\Ext{\mathop{\mathrm{Ext}}\nolimits}
\def\SheafHom{\mathop{\mathcal{H}\!\mathit{om}}\nolimits}
\def\SheafExt{\mathop{\mathcal{E}\!\mathit{xt}}\nolimits}
\def\Sch{\mathit{Sch}}
\def\Mor{\mathop{\mathrm{Mor}}\nolimits}
\def\Ob{\mathop{\mathrm{Ob}}\nolimits}
\def\Sh{\mathop{\mathit{Sh}}\nolimits}
\def\NL{\mathop{N\!L}\nolimits}
\def\CH{\mathop{\mathrm{CH}}\nolimits}
\def\proetale{{pro\text{-}\acute{e}tale}}
\def\etale{{\acute{e}tale}}
\def\QCoh{\mathit{QCoh}}
\def\Ker{\mathop{\mathrm{Ker}}}
\def\Im{\mathop{\mathrm{Im}}}
\def\Coker{\mathop{\mathrm{Coker}}}
\def\Coim{\mathop{\mathrm{Coim}}}

% Boxtimes
%
\DeclareMathSymbol{\boxtimes}{\mathbin}{AMSa}{"02}

%
% Macros for moduli stacks/spaces
%
\def\QCohstack{\mathcal{QC}\!\mathit{oh}}
\def\Cohstack{\mathcal{C}\!\mathit{oh}}
\def\Spacesstack{\mathcal{S}\!\mathit{paces}}
\def\Quotfunctor{\mathrm{Quot}}
\def\Hilbfunctor{\mathrm{Hilb}}
\def\Curvesstack{\mathcal{C}\!\mathit{urves}}
\def\Polarizedstack{\mathcal{P}\!\mathit{olarized}}
\def\Complexesstack{\mathcal{C}\!\mathit{omplexes}}
% \Pic is the operator that assigns to X its picard group, usage \Pic(X)
% \Picardstack_{X/B} denotes the Picard stack of X over B
% \Picardfunctor_{X/B} denotes the Picard functor of X over B
\def\Pic{\mathop{\mathrm{Pic}}\nolimits}
\def\Picardstack{\mathcal{P}\!\mathit{ic}}
\def\Picardfunctor{\mathrm{Pic}}
\def\Deformationcategory{\mathcal{D}\!\mathit{ef}}
\usepackage{cleveref}
\crefname{originallemma}{original lemma}{original lemmas}
\Crefname{originallemma}{Original lemma}{Original lemmas}
\makeatletter
\newlabel{more-algebra-lemma-extended-alternating-Cech-is-colimit-koszul}{{0913}{source}{}{sourcepointer-1}{}}
\newlabel{more-algebra-lemma-extended-alternating-Cech-is-colimit-koszul@cref}{{[originallemma][0][]0913}{[1][0][]source}{}{}{}}
\makeatother
\crefname{originallemma}{original lemma}{original lemmas}
\Crefname{originallemma}{Original lemma}{Original lemmas}
\makeatletter
\newlabel{definition-koszul-complex}{{0623}{source}{}{sourcepointer-2}{}}
\newlabel{definition-koszul-complex@cref}{{[originallemma][0][]0623}{[1][0][]source}{}{}{}}
\makeatother
\crefname{originallemma}{original lemma}{original lemmas}
\Crefname{originallemma}{Original lemma}{Original lemmas}
\makeatletter
\newlabel{section-koszul}{{0621}{source}{}{sourcepointer-3}{}}
\newlabel{section-koszul@cref}{{[originallemma][0][]0621}{[1][0][]source}{}{}{}}
\makeatother
\title{The extended alternating Cech complex as a colimit of Koszul complexes}
\author{The Stacks Project Authors}
\date{stacks-project:more-algebra.tex:local-snapshot-20261012; GNU FDL 1.2 or later}
\begin{document}
\maketitle
\setcounter{subsection}{1}
\noindent\textit{Editorial note.} All delivered mathematical text is the author's own, copied from the pinned Stacks Project chapter file more-algebra.tex; no statement or proof step was rewritten, reordered or removed.\par\medskip
\noindent\textit{Editorial note.} The chapter file pinned-original-source.tex carries only \textbackslash{}input\{preamble\} on its first line, so no preamble text lives inside the chapter. The retained file sources/stacks-more-algebra/retained-preamble.tex is a mechanical, line-for-line copy of the project's single shared preamble as staged in the snapshot (Xy-pic setup plus theorem declarations and mathematical macros). Nothing was added, removed or redefined, and the excerpt serves only as preamble, never as a span coordinate.\par\medskip
\noindent\textit{Editorial note.} Both counted spans start at the \textbackslash{}begin line of their environment, so the statement and the author's proof are delivered whole; surrounding prose of the section is deliberately outside the counted spans and is carried as standard prerequisites instead.\par\medskip
\noindent\textit{Editorial note.} The citation fields (source\_id, source\_version, authors, source\_url, license.name, source\_title) are identical to those of every other card this lane has produced from this chapter, so the batch selector treats them as one source.\par\medskip
\noindent\textit{Allowed standard prerequisites (editorial).}\begin{itemize}
\item Koszul complex of a finite sequence and its differential on exterior powers. constructed in the same chapter, source lines 6620-6700 and 7100-7298. Supplies the complexes whose colimit is being compared.
\item Extended alternating Cech complex, defined through the localisations at products of the elements. defined in the same section, source lines 6908-6944, immediately before the counted statement. Supplies the target of the comparison.
\item Colimits of complexes are computed degreewise and a map out of a colimit is a compatible system of maps on the finite pieces. used at proof lines 7120-7190. Makes the comparison maps and their compatibility meaningful.
\end{itemize}\medskip
\ifdefined\hypertarget\hypertarget{sourcepointer-1}{}\fi
\noindent\textit{Source reference.} Original lemma 0913: focused-seeds/AGP/sources/stacks-snapshot-20261011/tags.txt, source lines 11431--11431. The lemma's own canonical tag, quoted so that readers can resolve the label inside the delivered text; it points at the counted statement itself and carries no proof content.\par\medskip
\ifdefined\hypertarget\hypertarget{sourcepointer-2}{}\fi
\noindent\textit{Source reference.} Original lemma 0623: focused-seeds/AGP/sources/stacks-snapshot-20261011/tags.txt, source lines 7782--7782. Navigation pointer to a statement the counted proof refers to but does not contain; nothing load-bearing is delegated to it.\par\medskip
\ifdefined\hypertarget\hypertarget{sourcepointer-3}{}\fi
\noindent\textit{Source reference.} Original lemma 0621: focused-seeds/AGP/sources/stacks-snapshot-20261011/tags.txt, source lines 7780--7780. Navigation pointer to a statement the counted proof refers to but does not contain; nothing load-bearing is delegated to it.\par\medskip
\begin{lemma}
\label{lemma-extended-alternating-Cech-is-colimit-koszul}
Let $R$ be a ring. Let $f_1, \ldots, f_r \in R$. The extended alternating
{\v C}ech complex
$$
R \to \bigoplus\nolimits_{i_0} R_{f_{i_0}} \to
\bigoplus\nolimits_{i_0 < i_1} R_{f_{i_0}f_{i_1}} \to
\ldots \to R_{f_1\ldots f_r}
$$
is a colimit of the Koszul complexes $K(R, f_1^n, \ldots, f_r^n)$; see
proof for a precise statement.
\end{lemma}

\begin{proof}
We urge the reader to prove this for themselves.
Denote $K(R, f_1^n, \ldots, f_r^n)$ the Koszul complex of
Definition \ref{definition-koszul-complex} viewed as a cochain complex
sitting in degrees $0, \ldots, r$. Thus we have
$$
K(R, f_1^n, \ldots, f_r^n) :
0 \to \wedge^r(R^{\oplus r}) \to
\wedge^{r - 1}(R^{\oplus r}) \to \ldots \to
R^{\oplus r} \to R \to 0
$$
with the term $\wedge^r(R^{\oplus r})$ sitting in degree $0$.
Let $e^n_1, \ldots, e^n_r$ be the standard basis of $R^{\oplus r}$.
Then the elements $e^n_{j_1} \wedge \ldots \wedge e^n_{j_{r - p}}$ for
$1 \leq j_1 < \ldots < j_{r - p} \leq r$ form a basis for the term in
degree $p$ of the Koszul complex. Further, observe that
$$
d(e^n_{j_1} \wedge \ldots \wedge e^n_{j_{r - p}}) =
\sum (-1)^{a + 1} f_{j_a}^n e^n_{j_1} \wedge \ldots \wedge \hat e^n_{j_a} \wedge
\ldots \wedge e^n_{j_{r - p}}
$$
by our construction of the Koszul complex in Section \ref{section-koszul}.
The transition maps of our system
$$
K(R, f_1^n, \ldots, f_r^n) \to K(R, f_1^{n + 1}, \ldots, f_r^{n + 1})
$$
are given by the rule
$$
e^n_{j_1} \wedge \ldots \wedge e^n_{j_{r - p}}
\longmapsto
f_{i_0} \ldots f_{i_{p - 1}}
e^{n + 1}_{j_1} \wedge \ldots \wedge e^{n + 1}_{j_{r - p}}
$$
where the indices $1 \leq i_0 < \ldots < i_{p - 1} \leq r$ are such that
$\{1, \ldots r\} =
\{i_0, \ldots, i_{p - 1}\} \amalg \{j_1, \ldots, j_{r - p}\}$.
We omit the short computation that shows this is compatible with differentials.
Observe that the transition maps are always $1$ in degree $0$
and equal to $f_1 \ldots f_r$ in degree $r$.

\medskip\noindent
Denote $K^p(R, f_1^n, \ldots, f_r^n)$ the term of degree $p$ in the Koszul
complex. Observe that for any $f \in R$ we have
$$
R_f = \colim (R \xrightarrow{f} R \xrightarrow{f} R \to \ldots )
$$
Hence we see that in degree $p$ we obtain
$$
\colim K^p(R, f_1^n, \ldots f_r^n) =
\bigoplus\nolimits_{1 \leq i_0 < \ldots < i_{p - 1} \leq r}
R_{f_{i_0} \ldots f_{i_{p - 1}}}
$$
Here the element $e^n_{j_1} \wedge \ldots \wedge e^n_{j_{r - p}}$
of the Koszul complex above maps in the colimit
to the element $(f_{i_0} \ldots f_{i_{p - 1}})^{-n}$ in the
summand $R_{f_{i_0} \ldots f_{i_{p - 1}}}$
where the indices are chosen such that
$\{1, \ldots r\} = \{i_0, \ldots, i_{p - 1}\}
\amalg \{j_1, \ldots, j_{r - p}\}$.
Thus the differential on this complex is given by
$$
d(1\text{ in }R_{f_{i_0} \ldots f_{i_{p - 1}}}) =
\sum\nolimits_{i \not \in \{i_0, \ldots, i_{p - 1}\}}
(-1)^{i - t}\text{ in }
R_{f_{i_0} \ldots f_{i_t} f_i f_{i_{t + 1}} \ldots f_{i_{p - 1}}}
$$
Thus if we consider the map of complexes given in degree $p$
by the map
$$
\bigoplus\nolimits_{1 \leq i_0 < \ldots < i_{p - 1} \leq r}
R_{f_{i_0} \ldots f_{i_{p - 1}}}
\longrightarrow
\bigoplus\nolimits_{1 \leq i_0 < \ldots < i_{p - 1} \leq r}
R_{f_{i_0} \ldots f_{i_{p - 1}}}
$$
determined by the rule
$$
1\text{ in }R_{f_{i_0} \ldots f_{i_{p - 1}}}
\longmapsto
(-1)^{i_0 + \ldots + i_{p - 1} + p}\text{ in }R_{f_{i_0} \ldots f_{i_{p - 1}}}
$$
then we get an isomorphism of complexes from
$\colim K(R, f_1^n, \ldots, f_r^n)$ to the
extended alternating {\v C}ech complex defined in this section.
We omit the verification that the signs work out.
\end{proof}

\end{document}
